\documentclass[tikz, border=8pt]{standalone}
\usepackage{tikz}
\usetikzlibrary{calc}
\definecolor{navy}{HTML}{14213D}
\definecolor{teal}{HTML}{0E7C7B}
\definecolor{emas}{HTML}{C9A227}
\definecolor{abu}{HTML}{E8E8E8}

% Sudut refleks: sudut yang besarnya lebih dari 180 derajat
% tetapi kurang dari 360 derajat.
%
% Gambar 1: dua sinar dari titik O. Sinar OA mendatar (0 derajat),
% sinar OB ke arah 240 derajat. Sudut kecil 120 derajat (busur emas,
% lewat bawah) dan sudut refleks 240 derajat (busur teal, lewat atas).
% Jumlah keduanya 360 derajat.
% Gambar 2: jam dinding pukul 04.00. Sudut kecil antara kedua jarum
% 4 x 30 = 120 derajat; sudut refleksnya 360 - 120 = 240 derajat.

\begin{document}

% ---------------- Gambar 1: sudut refleks 240 derajat ----------------
\begin{tikzpicture}[x=1.3cm, y=1.3cm]
  % Sinar OA (0 derajat) dan OB (240 derajat)
  \draw[navy, very thick] (0,0) -- (2.6,0);
  \draw[navy, very thick] (0,0) -- (-1.3,-2.25);
  % Busur refleks: dari 0 derajat sampai 240 derajat
  % (TikZ menggambar busur berlawanan arah jarum jam)
  \draw[teal, very thick] (0,0) ++(0:1.5) arc (0:240);
  \node[anchor=east, font=\small\bfseries, teal] at (-1.7,0.95) {$240^\circ$};
  % Busur kecil: dari -120 derajat sampai 0 derajat (lewat bawah)
  \draw[emas, very thick] (0,0) ++(-120:1.0) arc (-120:0);
  \node[anchor=west, font=\small\bfseries, emas] at (0.75,-1.2) {$120^\circ$};
  % Titik O, A, B
  \fill[navy] (0,0) circle (0.06cm);
  \fill[navy] (2.6,0) circle (0.06cm);
  \fill[navy] (-1.3,-2.25) circle (0.06cm);
  \node[anchor=west, font=\small\bfseries, navy] at (0.15,0.18) {O};
  \node[anchor=south, font=\small\bfseries, navy] at (2.6,0.12) {A};
  \node[anchor=east, font=\small\bfseries, navy] at (-1.45,-2.25) {B};
\end{tikzpicture}

% ---------------- Gambar 2: jam dinding pukul 04.00 ----------------
\begin{tikzpicture}[x=1.1cm, y=1.1cm]
  % Lingkaran jam
  \draw[navy, very thick] (0,0) circle (2.2);
  % Dua belas tanda jam
  \foreach \a in {0,30,...,330} {
    \draw[navy, thick] (\a:2.05) -- (\a:2.2);
  }
  % Angka 12, 3, 6, 9
  \node[font=\small\bfseries, navy] at (0,1.75) {12};
  \node[font=\small\bfseries, navy] at (1.75,0) {3};
  \node[font=\small\bfseries, navy] at (0,-1.75) {6};
  \node[font=\small\bfseries, navy] at (-1.75,0) {9};
  % Busur refleks: dari 90 derajat sampai 330 derajat (lewat kiri dan bawah)
  \draw[teal, very thick] (0,0) ++(90:1.45) arc (90:330);
  \node[anchor=east, font=\small\bfseries, teal] at (-1.5,-0.85) {$240^\circ$};
  % Busur kecil: dari -30 derajat sampai 90 derajat (lewat kanan)
  \draw[emas, very thick] (0,0) ++(-30:1.1) arc (-30:90);
  \node[anchor=west, font=\small\bfseries, emas] at (1.35,0.78) {$120^\circ$};
  % Jarum menit ke 12 dan jarum jam ke 4 (sudut -30 derajat)
  \draw[navy, very thick] (0,0) -- (0,1.7);
  \draw[navy, very thick] (0,0) -- (1.47,-0.85);
  \fill[navy] (0,0) circle (0.08cm);
\end{tikzpicture}

\end{document}
